% se puede agregar la opción [english] para 
%  memorias o tesis en inglés (borrando el archivo .aux)
\documentclass{umemoria} 

\depto{Departamento de Ciencias de la Computación}
\author{Dylan Nicolás Riquelme Vergara}
\title{Diseño y desarrollo de sistema de evaluación automático de ejercicios de seguridad ofensiva}

% incluir ambos comandos para una doble titulación
%  o quitar el comando que no aplica
\memoria{Ingeniero Civil en Computación}
%\tesis{Doctor en ???} % incluir solo este comando para doctorados

% puede haber varios profesores guía seperados por coma;
% pero si es una memoria, solo puede haber un profesor guía
\guia{Eduardo Riveros Roca} 

% puede haber varios profesores co-guía seperados por coma;
% pero si es una memoria, el profesor co-guía será el primer
% integrante de la comisión
%\coguia{Nombre Completo Co-Guía} % incluir en caso de co-guía de *tesis*

%\cotutela{Nombre Institución} % incluir en caso de cotutela

\comision{Nombre Completo Uno,Nombre Completo Dos,Nombre Completo Tres}

%\auspicio{Nombre Institución} % incluir en caso de recibir financiamiento

% tiene que ser el año en que se da el examen de título/grado (defensa)
%\anho{2021} % incluir solo para reemplazar el año actual

\usepackage{lipsum}
\usepackage{xcolor}

\begin{document}

\frontmatter
\maketitle

\begin{resumen}
    \textbf{\colorbox{yellow}{PENDIENTE PENDIENTE PENDIENTE}}
\end{resumen}

% opcional: incluir para tesis en inglés;
%  en este caso hay que tener el resumen y abstract
%   en ambos idiomas
%\begin{abstract}
%\lipsum[1-4]
%\end{abstract}

\begin{dedicatoria}
\textbf{\colorbox{yellow}{PENDIENTE PENDIENTE PENDIENTE}}
\end{dedicatoria}

\begin{thanks}
    \begin{itemize}
        \item Cami
        \item Familia
        \item Amigos Colegio
        \item Amigos Universidad
        \item Eduardo
        \item CLCERT
        \item DCC
        \item Maribel
        \item 
    \end{itemize}
\end{thanks}

\tableofcontents
\listoftables % opcional
\listoffigures % opcional

\mainmatter

\chapter{Introducción}
\label{sec:introduccion}

Actualmente, la seguridad en el desarrollo de software es fundamental. Sin
embargo, las malas prácticas y la falta de conocimiento siguen existiendo a
pesar del aumento de amenazas y la mejora continua de capacidades de los
atacantes.

Existe abundante literatura respecto a las vulnerabilidades con las que se puede
encontrar un desarrollador, como el OWASP Top 10~\cite{owasptop10}, que cataloga
las vulnerabilidades de software más críticas, incluyendo inyección, pérdida de
control de acceso, fallas criptográficas, entre otras. Sin embargo, la mejor
forma de conocer el riesgo real de las vulnerabilidades es entendiendo cómo
funcionan, para lo cual un enfoque práctico permite validar su correcto
entendimiento. Las investigaciones demuestran que el aprendizaje práctico es
significativamente más efectivo, con estudios mostrando que los participantes
retienen el 75\% de la información a través de práctica \textit{hands-on},
comparado con solo el 5\% mediante métodos teóricos
tradicionales~\cite{sanshandson}.

El problema a abordar en este trabajo es facilitar, tanto al equipo docente como
a los estudiantes, confirmar sus conocimientos teóricos a partir de una
plataforma automatizada que valide sus avances. Las y los desarrolladores
necesitan entornos controlados y seguros donde puedan experimentar y explorar
vulnerabilidades posibles dentro de una aplicación sin que esta sea de
producción. Para ello, se diseñó y desarrolló un sistema de evaluación
automático de ejercicios de ciberseguridad ofensiva.

Este proyecto se aproxima al problema desde tres perspectivas complementarias.
Desde lo \textbf{pedagógico}, se requiere una forma de evaluar el progreso que
capture la trayectoria completa del aprendizaje, no solo el resultado final. La
validación automática basada en el estado del sistema permite registrar cada
paso del estudiante durante el ataque, proporcionando trazabilidad real del
proceso de aprendizaje sin recurrir a soluciones artificiales como múltiples
\textit{flags} intermedias. Desde lo \textbf{técnico}, existe una oportunidad de
contribuir al ecosistema de herramientas \textit{open source} de detección de
intrusiones, publicando código, patrones y metodologías que permitan que la
comunidad académica evalúe y mejore estos enfoques. Desde lo \textbf{personal},
este proyecto representa la oportunidad de desarrollar experiencia práctica en
el diseño e implementación de sistemas complejos con múltiples componentes,
arquitecturas de capas bien definidas, y desafíos de observabilidad cercanos a
los que enfrenta la industria de ciberseguridad.

\section{Objetivos}

\subsection{Objetivo General}

Desarrollar una plataforma educativa autohospedable que permita crear entornos
de práctica de seguridad ofensiva personalizables, con validación automática del
progreso basada en estado de los sistemas comprometidos, facilitando el
aprendizaje práctico de vulnerabilidades sin depender de validaciones manuales.

\subsection{Objetivos Específicos}

\begin{enumerate}
    \item Diseñar e implementar una arquitectura modular que permita desplegar y
    gestionar múltiples entornos vulnerables simultáneamente mediante
    tecnologías de virtualización o contenerización, asegurando aislamiento y
    facilidad de configuración.
    
    \item Desarrollar un sistema de configuración declarativa que permita
    definir las condiciones de éxito de un ejercicio mediante archivos YAML,
    de modo que crear un ejercicio con validación automática requiera
    describir qué debe observarse y no programar la lógica de detección ni
    modificar la plataforma.

    \item Implementar el motor de validación automática que evalúe el estado de
    los sistemas comprometidos, determinando el progreso sin depender de
    \textit{flags} predefinidas, ya sea por archivos modificados, servicios
    comprometidos, usuarios creados, entre otros.

    \item Crear herramientas de seguimiento del progreso, acceso a enunciados de
    ejercicios y administración del sistema para educadores y estudiantes,
    implementadas inicialmente mediante línea de comandos.

    \item Desplegar la plataforma en la infraestructura del laboratorio de
    ciberseguridad CLCERT y verificar allí su operación completa: instalación,
    ejecución de extremo a extremo de los escenarios y acceso remoto de un
    participante.\footnote{Este objetivo, formulado originalmente como ``poner
    la plataforma a disposición del CLCERT para su evaluación en una instancia
    real de uso'', se precisa aquí con un criterio verificable.}
\end{enumerate}

\section{Evaluación}

La forma de evaluar este trabajo fue verificando la capacidad del sistema para
generar diferentes tipos de ejercicios con validación automática. El trabajo se
consideró exitoso si era posible crear ejercicios de seguridad ofensiva con
validación automática para al menos tres categorías entre: aplicaciones y
servicios web, bases de datos, \textit{pwning}, análisis forense e ingeniería
reversa. Se entendió por validación automática aquella que marca el ejercicio
como completado automáticamente cuando el participante logra explotar una
vulnerabilidad, acceder a un archivo específico, comprometer un servicio u
obtener privilegios de usuario, sin requerir la entrada manual de \textit{flags}
o códigos predefinidos.

\section{Solución Propuesta}

La solución consiste en una plataforma que permite a los educadores cargar
imágenes de sistemas vulnerables y definir reglas de validación mediante
archivos YAML, generando automáticamente entornos de práctica personalizados con
validación automática.

\subsection{Arquitectura General}

El sistema implementa una arquitectura dividida en cuatro capas con
responsabilidades claramente delimitadas. Esta arquitectura emergió del análisis
de pruebas de concepto~\cite{flexipwn-poc} desarrolladas durante el trabajo
previo al semestre, donde se identificó que la separación de responsabilidades
entre el sistema vulnerable, el mecanismo de detección y la lógica de evaluación
es fundamental para lograr flexibilidad y extensibilidad.

\begin{description}
    \item[Capa 1 — Sistemas aislados] Contiene todos los sistemas y redes que
    componen el escenario de ataque. La implementación actual los materializa
    como contenedores, aunque la arquitectura es agnóstica al origen de
    aislamiento y podría sostenerse a futuro sobre tecnologías de
    virtualización. Se divide en dos subcapas: la Capa 1a, el sistema vulnerable
    objetivo (el contenedor vulnerable), y la Capa 1b, el entorno atacante desde
    el cual el estudiante ejecuta el ataque (el contenedor atacante).

    \item[Capa 2 — Monitoreo pasivo] Observa pasivamente el estado de los
    contenedores con acceso de lectura externo al sistema. La única excepción a
    esa lectura externa pura es la captura de tráfico de red, que requiere un
    contenedor auxiliar (\textit{sidecar}) corriendo \texttt{tcpdump}, por ser
    la forma adecuada de capturar el tráfico del escenario. Los escenarios de
    referencia asumen tráfico en claro, sin TLS, para que la captura resulte
    legible. La responsabilidad de esta capa es solo generar un registro
    estructurado de eventos detectables: modificaciones del sistema de archivos,
    procesos creados, entradas de \textit{log} y tráfico o conexiones de red,
    entre
    otros. Su comportamiento es análogo al de un SIEM (Security Information and
    Event Management) acotado a un único escenario de ataque.

    \item[Capa 3 — Motor de evaluación] Consume los eventos generados por la
    Capa 2 y evalúa reglas declarativas definidas en archivos YAML que
    especifican las condiciones de victoria. Determina automáticamente cuándo un
    participante ha completado objetivos sin requerir la entrada manual de
    \textit{flags}.

    \item[Capa 4 — Administración] Gestiona el ciclo de vida completo de los
    entornos de práctica y mantiene registro del progreso de todos los
    participantes. Proporciona interfaces de línea de comandos principalmente
    para educadores, incluyendo vistas agregadas para seguir varios entornos
    en paralelo.
\end{description}

\subsection{Flujo principal}

\begin{enumerate}
    \item \textbf{Carga de escenarios}: el educador entrega una imagen del
    sistema vulnerable junto con un archivo YAML que define los objetivos del
    ejercicio, las condiciones de validación y los criterios de éxito.

    \item \textbf{Generación de entornos}: la plataforma instancia la imagen
    proporcionada, creando un entorno aislado por cada participante que
    intentará explotar la vulnerabilidad.

    \item \textbf{Validación automática}: el motor de validación examina
    constantemente el estado del contenedor, comparando el estado actual contra
    las reglas definidas en el archivo YAML.

    \item \textbf{Gestión}: tanto educadores como estudiantes utilizan
    herramientas de línea de comandos para interactuar con sus respectivos
    entornos y monitorear el avance.
\end{enumerate}

\subsection{Perfil de usuario esperado}
\label{sec:perfil-usuario}

FlexiPwn asume un educador con manejo de línea de comandos y familiaridad
básica con contenedores: debe poder instalar Docker en modo \textit{rootless}
una vez, construir o adaptar una imagen del sistema vulnerable y escribir un
archivo YAML. No es el tipo de usuario final que solo tiene que saber usar una aplicación web, sino
que corresponde al perfil habitual de quien dicta un curso de
seguridad ofensiva. Lo que la plataforma elimina es la necesidad de programar
el mecanismo de detección, es decir, configurar mecanismos para observar el
sistema objetivo, capturar y correlacionar eventos, y escribir el código que decide cuándo un objetivo se
cumple, que es la parte costosa de montar validación automática. Ese es el
sentido en que la plataforma resulta accesible: no porque rebaje los
requisitos de infraestructura, sino porque sustituye el trabajo de
instrumentación por una descripción declarativa.

Además, ese piso técnico podría considerarse el más bajo entre las alternativas
autohospedables analizadas en el Capítulo~\ref{sec:estado-del-arte}: CTFd
requiere desplegar un servicio web con base de datos, y los \textit{cyber
ranges} como Tectonic o CyberRangeCZ se operan con herramientas de
infraestructura como código, entre ellas Packer, Terraform y Ansible, sobre
proveedores de nube o hipervisores.

\section{Declaración de uso de inteligencia artificial}

En la elaboración de esta memoria se utilizaron herramientas de inteligencia
artificial generativa, concretamente los modelos Claude de Anthropic (familias
Opus y Sonnet 4.x), empleados principalmente a través de la aplicación de
escritorio en su modo Cowork. Estas herramientas se usaron a lo largo de todo
el trabajo: en la revisión de literatura y la comparación de
plataformas, en la discusión de decisiones de arquitectura, 
en la generación y refactorización de código, en la redacción y
edición de este documento,  y en tareas de apoyo
como el control de versiones y la organización del proyecto. Todo el contenido
generado o asistido fue revisado, evaluado y validado por el autor,
quien asume plena responsabilidad sobre el trabajo presentado. La
inteligencia artificial se empleó como herramienta de apoyo y no se reconoce
como coautoría.

El modo Cowork se utilizó como un acelerador del trabajo: permitió delegar la
realización del trabajo, como escribir código, redactar y editar, ejecutar y
verificar, conservando el rol del autor como arquitecto de las decisiones.
Trabajar de esta forma permitió mantener la constancia durante todo el semestre,
compatibilizando con las labores del autor como desarrollador de aplicaciones
móviles. La integración de la herramienta estuvo acompañado de la verificación:
la relectura crítica de cada resultado, la comparación de las afirmaciones y
planeamientos contra el código fuente, y el uso de una memoria persistente entre
sesiones para mantener el contexto y evitar contradicciones. El criterio
principal fue no dar por válido ningún resultado por el solo hecho de que
``funciona'' a primera vista.

Los prompts siguieron unos patrones recurrentes. Para una tarea de inicio a fin,
se partía entregando el contexto (usualmente un resumen de lo trabajado en
sesiones anteriores), seguido de lo que había que hacer y de los puntos que
valía la pena discutir, pidiendo siempre alternativas para elegir la
implementación más adecuada. Se cerraba con algún criterio de validación y con
notas que precisaban la ejecución. A mitad de una tarea el patrón cambiaba: se
revisaba lo ejecutado y se respondía con un punteo extenso de dudas, cambios
deseados y aspectos por considerar, y solo se confirmaba el avance al paso
siguiente cuando el resultado era satisfactorio. La mayoría de veces, se
recordaba al modelo lo que debería tener en memoria, para asegurar un correcto
alineamiento de la tarea a realizar. Los prompts breves se usaban para algunos
casos puntuales cuya discusión convenía expandir, y al final de una sesión
cargada se solicitaban los cierres habituales: formateo, guardado en memoria,
mensaje de \textit{commit} si se requería y un prompt de traspaso que orientara
a la sesión siguiente.

A modo ilustrativo, dos ejemplos representativos del tipo de instrucción
entregada:

\begin{itemize}
    \item Un prompt que funcionó bien: ``Retoma el informe: primero una foto del
    estado del repositorio con los commits nuevos. Luego valida afirmación por
    afirmación contra el código, marcando cada una como verdadera, falsa o
    imprecisa con su archivo y línea. Entrega un reporte ordenado por criticidad
    antes de editar y espera mi confirmación. Reglas: 80 caracteres, prosa en
    presente y opciones ante cada alternativa.'' Funcionó porque aportaba
    contexto, descomponía la tarea, exigía validación y opciones, y fijaba
    restricciones explícitas.
    \item Un prompt que funcionó mal: una instrucción amplia y sin contexto como
    ``mejora el capítulo de implementación'', que al no acotar el objetivo ni
    los criterios de aceptación producía cambios genéricos o de alcance excesivo
    que luego exigían una corrección considerable. La lección es que los prompts
    pobres en contexto y sin criterios claros rendían resultados de bajo valor.
\end{itemize}

Surgió un momento interesante en la etapa final del trabajo, al estudiar una
prueba de concepto de escape de contenedores para comprender el funcionamiento
de este tipo de \textit{exploit} y respaldar el análisis de contención del
entorno. Al
tratar un tema sensible, las primeras consultas sobre una prueba de concepto
realista del CVE~\cite{cve-runc-2025, poc-cve-2025-31133} en cuestión fueron
rechazadas por las políticas de uso de la herramienta. Tras solicitar al equipo
de Anthropic un acceso ampliado y explicar su propósito educativo dentro de este
trabajo de ciberseguridad, la consulta pudo retomarse. A partir de ahí, el
modelo resultó de gran ayuda para seguir el \textit{exploit} paso a paso,
consiguiendo
compararlo con otros CVE relacionados, distinguiendo, por ejemplo, los que se
ejecutan desde el interior del contenedor de los que afectan directamente al
\textit{host}.

Si el trabajo se reiniciara, se cambiarían algunas formas del uso de la
inteligencia artificial. Se establecería antes una rutina de revisión y
verificación sistemática, para no depender de que un resultado ``parezca''
correcto. Además, se sería más disciplinado con la memoria persistente y los
traspasos entre sesiones, que resultaron clave para mantener la coherencia a lo
largo de un trabajo asistido por IA.

\chapter{Estado del Arte}

\section{Plataformas Comerciales}
\section{Plataformas con Validación Automática}
\section{Plataformas de Estáticas}
\section{Plataformas de Tipo CTF}
\section{Plataformas Avanzadas}
\section{Oportunidad Identificada}
\chapter{Diseño}
\label{sec:diseño}

Este capítulo describe las decisiones de diseño del sistema FlexiPwn, incluyendo
la arquitectura de capas, los contratos entre ellas, el modelo de datos y el
modelo de aislamiento de entornos.

\section{Arquitectura de Capas}

\subsection{Justificación}

La decisión de estructurar el sistema en capas responde a principios
arquitectónicos que se identificaron como necesarios a partir de las pruebas de
concepto desarrolladas previo al semestre. Las limitaciones encontradas
convergieron hacia un patrón común: la mezcla de responsabilidades entre el
sistema a vulnerar, el mecanismo de detección y la lógica de evaluación
resultaba en implementaciones frágiles, poco extensibles y fuertemente acopladas
a tecnologías específicas. A partir de este análisis emergieron cuatro
principios que fundamentan la arquitectura de cuatro capas.

\subsubsection*{Separación de responsabilidades}

Cada capa tiene una función específica y bien definida, lo que facilita el
desarrollo, \textit{testing} y mantenimiento independiente de cada componente.
Esta separación permite trabajar en capas distintas sin interferencia, y
simplifica la depuración cuando surgen problemas. En las pruebas de concepto
previas, la generación de eventos y la evaluación de condiciones de éxito
convivían en el mismo módulo, lo que hacía que agregar un nuevo tipo de fuente
de eventos o cambiar la lógica de evaluación requiriera modificar el mismo
archivo, aumentando el acoplamiento y la fragilidad del sistema. En FlexiPwn,
estas responsabilidades están explícitamente separadas: la Capa 2 solo observa y
genera eventos, la Capa 3 solo evalúa condiciones, y ninguna de las dos modifica
el estado de los entornos.

\subsubsection*{Reutilización y flexibilidad}

Al abstraer las capas inferiores mediante interfaces bien definidas, el sistema
puede soportar múltiples tecnologías de virtualización sin modificar las capas
superiores de validación y administración. Esto permite adaptar el sistema a
diferentes contextos e infraestructuras. La interfaz abstracta
\texttt{EnvironmentProvider} encapsula toda la complejidad necesaria para
distintas tecnologías de virtualización, como Docker \textit{rootless}: cambiar
la tecnología de virtualización solo requiere implementar esa interfaz en una
nueva clase, sin tocar el motor de evaluación ni la \textit{CLI} de
administración. Por ejemplo, una futura implementación \texttt{PodmanProvider} o
\texttt{KVMProvider} no tendría impacto en las capas 3 y 4.

\subsubsection*{Extensibilidad}

Nuevos tipos de validaciones o métodos de monitoreo pueden agregarse a las capas
2 y 3 sin modificar los sistemas vulnerables ni la administración. Por ejemplo,
agregar soporte para detectar exfiltración de datos mediante DNS
\textit{tunneling} solo requiere modificar la Capa 2 para capturar ese tipo de
eventos y la Capa 3 para evaluar las reglas correspondientes. El diseño de
registro de evaluadores mediante un patrón de \textit{registry} en la Capa 3
permite añadir nuevos tipos de \textit{target} sin modificar el motor de
evaluación central.

Esta extensibilidad opera principalmente sobre las capas 2 y 3, pero algunos
métodos de extracción de información requieren ampliar también la interfaz
expuesta por la Capa 1. La incorporación del monitor de red, por ejemplo, exigió
agregar el aprovisionamiento de un contenedor \textit{sidecar} de captura al
\texttt{EnvironmentProvider}; un futuro monitor basado en eBPF probablemente
requiera exponer un nuevo método para adjuntar programas al \textit{cgroup} del
contenedor vulnerable. La interfaz abstracta admite estos crecimientos siempre
que se respete el principio de pasividad: la Capa 1 expone capacidades de
observación, no las consume ni las interpreta.

\subsubsection*{Agnóstico a la implementación}

Los educadores que crean desafíos no necesitan conocer los detalles de cómo
funciona la detección; solo deben proporcionar una imagen vulnerable y un
archivo YAML con las reglas declarativas. Esto reduce significativamente la
barrera de entrada para crear contenido educativo. De igual forma, los sistemas
vulnerables de Capa 1 no saben que están siendo monitoreados, lo cual es un
principio crítico para la validez pedagógica del ejercicio: el entorno debe
comportarse exactamente como lo haría en producción, sin cooperar con el sistema
de detección. Este principio fue el aprendizaje más importante extraído de las
pruebas de concepto previas, donde la lógica de detección estaba parcialmente
implementada dentro del sistema vulnerable, violando la separación y limitando
la aplicabilidad del enfoque.

No obstante, ser agnóstico a la implementación no implica ser independiente del
tipo de aplicación. La eficacia de cada método de detección sigue dependiendo de
las características del sistema vulnerable expuesto. La detección de inyección
SQL \texttt{(log\_pattern + network\_payload)}, por ejemplo, asume que el motor
de base de datos elegido permite exponer un \textit{query log} legible por el
monitor (caso de MySQL con \texttt{general\_log}) y, cuando además se quiere
inspeccionar el payload SQL a nivel de red, que la conexión
\textit{cliente-servidor} opere sin TLS en el escenario, dado que el monitor de
red solo expone metadatos sobre tráfico cifrado. Un escenario que use una base
de datos que no exponga logs equivalentes obligaría al educador a definir
condiciones de éxito por otras señales. Esta dependencia es propia del dominio
del ejercicio y no del sistema de detección: el motor sigue ejecutando reglas
declarativas sin conocer detalles de la aplicación, pero la elección de qué
reglas son verificables sí depende de qué señales emite el sistema vulnerable.

\subsubsection*{\textit{Trade-offs} de la arquitectura}

Esta arquitectura introduce complejidad adicional en comparación con un diseño
monolítico, particularmente en la coordinación entre capas, la necesidad de
definir interfaces claras y el \textit{overhead} de comunicación entre
componentes. Sin embargo, los beneficios en términos de mantenibilidad,
portabilidad y capacidad de evolución del sistema justifican esta complejidad
inicial, especialmente considerando que se espera que el sistema crezca con
nuevos tipos de validadores, tecnologías de virtualización y escenarios de
ejercicios a futuro.

\subsection{Descripción de Capas}

\begin{figure}[H]
\centering \resizebox{\textwidth}{!}{%
\begin{tikzpicture}[
    % Estilos base
    every node/.style={font=\sffamily},
    arr/.style={-{Stealth[length=7pt]}, thick, gray!60!black},
    arrin/.style={-{Stealth[length=5pt]}, gray!50!black, semithick},
    % Sub-bloque interior
    sub/.style={rounded corners=3pt, draw=#1!30, fill=white,
                minimum height=1.2cm, text centered,
                font=\sffamily\large, inner sep=6pt},
    % Etiqueta de flecha
    lbl/.style={font=\sffamily\normalsize, text=gray!70!black, fill=white,
                inner sep=2pt},
    % Etiqueta vertical
    vlbl/.style={font=\sffamily\normalsize, text=gray!70!black, fill=white,
                 inner sep=2pt, rotate=90, anchor=south},
]

% ── Constantes de layout ────────────────────────────
\def\fullW{23}       % ancho capas completas
\def\halfW{10}      % ancho cada capa media
\def\midgap{2.8}     % espacio horizontal entre capas 3 y 2
\def\vgap{1.8}       % espacio vertical entre filas
\def\cx{9}           % centro X global

% ═══════════════════════════════════════════════════
% CAPA 4 — Administración (arriba, ancho completo)
% ═══════════════════════════════════════════════════
\node[rounded corners=5pt, draw=violet!40, fill=violet!6,
      minimum width=\fullW cm, minimum height=4.8cm,
      inner sep=12pt] (L4) at (\cx, 0) {};
\node[font=\sffamily\bfseries\Large, text=violet!70!black]
      at ([yshift=-20pt]L4.north) (t4) {Capa 4 — Administración};
% Sub-bloques
\node[sub=violet, text width=4.6cm, below=10pt of t4, xshift=-5.4cm]
      (c4a) {CLI\\[-2pt]{\normalsize Typer + Rich}};
\node[sub=violet, text width=4.6cm, right=20pt of c4a]
      (c4b) {Persistencia\\[-2pt]{\normalsize SQLModel + SQLite}};
\node[sub=violet, text width=4.6cm, right=20pt of c4b]
      (c4c) {Ciclo de vida\\[-2pt]{\normalsize Scenario, Run, User}};
% Flechas internas
\draw[arrin] (c4a.east) -- (c4b.west);
\draw[arrin] (c4c.west) -- (c4b.east);
% Bloque Modelos
\node[sub=violet, text width=15.2cm, below=10pt of c4b]
      (c4d) {Modelos: Scenario | Participant | ExerciseRun};

% ═══════════════════════════════════════════════════
% FILA MEDIA: Capa 3 (izquierda) + Capa 2 (derecha)
% ═══════════════════════════════════════════════════
\pgfmathsetmacro{\midY}{-4.4/2 - \vgap - 6.0/2}

% ── Capa 3 — Motor de evaluación ──────────────────
\pgfmathsetmacro{\layerIIIcx}{\cx - \midgap/2 - \halfW/2}
\node[rounded corners=5pt, draw=green!30!black!40, fill=green!5,
      minimum width=\halfW + 9.5cm, minimum height=6.0cm,
      inner sep=12pt, anchor=north]
      (L3) at (\layerIIIcx, \midY + 6.0/2) {};
\node[font=\sffamily\bfseries\Large, text=green!55!black]
      at ([yshift=-20pt]L3.north) (t3) {Capa 3 — Motor de evaluación};
% Schema YAML → Motor de reglas → Evaluador
\node[sub=green!70!black, text width=2cm, below=10pt of t3, xshift=-3.05cm]
      (c3a) {Schema YAML\\[-2pt]{\normalsize Pydantic v2}};
\node[sub=green!70!black, text width=2.4cm, right=12pt of c3a]
      (c3b) {Motor de reglas\\[-2pt]{\normalsize Python puro}};
\node[sub=green!70!black, text width=2.2cm, right=12pt of c3b]
      (c3c) {Evaluador\\[-2pt]{\normalsize any / all}};
% Flechas internas
\draw[arrin] (c3a.east) -- (c3b.west);
\draw[arrin] (c3b.east) -- (c3c.west);
% Tipos de target
\node[sub=green!70!black, text width=7cm, below=10pt of c3b]
      (c3d) {Tipos de target\\[-2pt]{\normalsize file\_created, process\_spawned, etc.}};

% ── Capa 2 — Monitoreo pasivo ─────────────────────
\pgfmathsetmacro{\layerIIcx}{\cx + \midgap/2 + \halfW/2}
\node[rounded corners=5pt, draw=orange!60!red!40, fill=orange!5!red!2,
      minimum width=\halfW cm, minimum height=6.0cm,
      inner sep=12pt, anchor=north]
      (L2) at (\layerIIcx, \midY + 6.0/2) {};
\node[font=\sffamily\bfseries\Large, text=orange!65!red]
      at ([yshift=-20pt]L2.north) (t2) {Capa 2 — Monitoreo pasivo};
% Los 3 monitores
\node[sub=orange!70!red, text width=2.2cm, below=28pt of t2, xshift=-2.4cm]
      (c2a) {Filesystem\\[-2pt]{\normalsize c.diff()}};
\node[sub=orange!70!red, text width=1.8cm, right=6pt of c2a]
      (c2b) {Procesos\\[-2pt]{\normalsize c.top()}};
\node[sub=orange!70!red, text width=1.8cm, right=6pt of c2b]
      (c2c) {Logs\\[-2pt]{\normalsize volúmenes}};
% Caja punteada orquestador
\node[rounded corners=3pt, draw=orange!50!red, densely dotted, thick,
      inner sep=8pt, fit=(c2a)(c2b)(c2c),
      label={[font=\sffamily\normalsize, text=orange!60!red]above:Monitor Orquestador}]
      (orch) {};
% Eventos JSON
\node[sub=orange!70!red, text width=7cm, below=10pt of orch]
      (c2d) {Eventos JSON estructurados\\[-2pt]{\normalsize timestamp, monitor\_type, etc.}};

% ═══════════════════════════════════════════════════
% CAPA 1 — Entornos virtualizados (abajo, ancho completo)
% ═══════════════════════════════════════════════════
\pgfmathsetmacro{\layerItop}{\midY - 6.0/2 - \vgap}
\node[rounded corners=5pt, draw=blue!40, fill=blue!5,
      minimum width=\fullW cm, minimum height=3.6cm,
      inner sep=12pt, anchor=north]
      (L1) at (\cx, \layerItop) {};
\node[font=\sffamily\bfseries\Large, text=blue!70!black]
      at ([yshift=-24pt]L1.north) (t1) {Capa 1 — Entornos virtualizados (Docker rootless)};
% 1a y 1b
\node[sub=blue, text width=7.5cm, below=10pt of t1, xshift=-4.8cm]
      (c1a) {\textbf{1a — Sistema vulnerable}\\[-2pt]
      {\normalsize Imagen Docker personalizada}};
\node[sub=blue, text width=7.5cm, right=2cm of c1a]
      (c1b) {\textbf{1b — Sistema atacante}\\[-2pt]
      {\normalsize Entorno y herramientas del estudiante}};
% Flecha Red
\draw[{Stealth[length=6pt]}-{Stealth[length=6pt]}, densely dashed,
      blue!60, thick] (c1a.east) -- (c1b.west)
      node[midway, fill=blue!5, font=\sffamily\large,
           text=blue!60!black, inner sep=3pt] {Red};

% ═══════════════════════════════════════════════════
% FLECHAS ENTRE CAPAS
% ═══════════════════════════════════════════════════

% Capa 4 ↔ Capa 3 (vertical)
\draw[arr] ([xshift=-16pt]L4.south -| L3.north)
    -- ([xshift=-16pt]L3.north)
    node[midway, lbl, left=2pt] {Config};
\draw[arr] ([xshift=16pt]L3.north)
    -- ([xshift=16pt]L4.south -| L3.north)
    node[midway, lbl, right=2pt] {Resultados};

% Capa 2 → Capa 3 (horizontal, eventos JSON)
\draw[arr] (L2.west |- {$(L2.north)!0.5!(L2.south)$})
    -- (L3.east |- {$(L3.north)!0.5!(L3.south)$})
    node[midway, lbl, above=5pt] {Eventos JSON};

% Capa 1 ↑ Capa 2 (3 flechas: filesystem, procesos, logs)
\draw[arr] ([xshift=-3cm]L1.north -| L2.south)
    -- ([xshift=-3cm]L2.south)
    node[midway, lbl, left=1pt] {Lectura fs};
\draw[arr] (L1.north -| L2.south)
    -- (L2.south)
    node[midway, lbl, left=1pt] {Lectura proc};
\draw[arr] ([xshift=3cm]L1.north -| L2.south)
    -- ([xshift=3cm]L2.south)
    node[midway, lbl, left=1pt] {Lectura logs};

% Capa 4 → Capa 1 (ciclo de vida, lateral derecho) — texto vertical
\draw[arr, rounded corners=6pt]
    (L4.east) -- ++(1.0, 0) |- (L1.east)
    node[pos=0.255, vlbl] {Ciclo de vida};

% ═══════════════════════════════════════════════════
% ACTORES (debajo de Capa 1)
% ═══════════════════════════════════════════════════
\node[below=1cm of L1.south, xshift=-3.5cm,
      font=\sffamily\large, align=center]
      (est) {\textbf{Estudiante}\\[-2pt]{\normalsize vía CLI o SSH}};
\node[below=1cm of L1.south, xshift=4.5cm,
      font=\sffamily\large, align=center]
      (edu) {\textbf{Educador}\\[-2pt]{\normalsize vía CLI}};

% Estudiante → Capa 1 (directo, vertical)
\draw[arr] (est.north) -- (est.north |- L1.south)
    node[midway, lbl, left=2pt] {Ataca 1a desde 1b};

% Educador → Capa 4 (rodea por la derecha) — texto vertical
\draw[arr, rounded corners=6pt]
    (edu.east) -- ++(\fullW/2 - 4.5 + 0.9, 0) |- ([yshift=-10pt]L4.east)
    node[pos=0.3, vlbl] {Administra};

% Educador → Schema YAML Capa 3 (rodea por la izquierda, pasa debajo de Estudiante) — texto vertical
\draw[arr, rounded corners=6pt]
    (edu.south) -- ++(0, -0.6)
    -| ([xshift=-0.7cm]L3.west |- est.south)
    |- (L3.west |- {$(L3.north)!0.35!(L3.south)$})
    node[pos=0.4, vlbl] {YAML escenario};

\end{tikzpicture}
%
}% fin resizebox
\caption{Arquitectura de cuatro capas de FlexiPwn.}
\label{fig:arquitectura}
\end{figure}

\subsubsection{Capa 1: Sistemas Virtualizados}

Contiene en estado virtualizado todos los sistemas y redes que componen el
escenario de ataque. Se divide en dos subcapas con roles diferenciados pero
complementarios.

La \textbf{Capa 1a} representa el sistema vulnerable objetivo. Esta subcapa
autocontiene todo el estado (archivos, memoria, procesos) necesario para
determinar si un ataque fue exitoso. Es accesible al participante pero no
comparte información de forma activa con otras capas del sistema, y no sabe que
está siendo monitoreada.

La \textbf{Capa 1b} proporciona al participante un entorno con herramientas para
ejecutar ataques. No tiene más acceso al sistema vulnerable que el necesario
para ejecutar el ataque, típicamente conexión a la misma red interna.

\subsubsection{Capa 2: Generación de Eventos}

Tiene acceso de lectura externo al estado de los contenedores de Capa 1. Su
única responsabilidad es generar un registro estructurado de eventos de
seguridad detectables: cambios en el sistema de archivos, nuevos procesos,
entradas de log de aplicación y tráfico de red observado. Su comportamiento es
cercano al de un SIEM dedicado al escenario de ataque.

El principio fundamental de esta capa es la \textbf{pasividad total}: no ejecuta
ningún comando dentro del contenedor vulnerable ni modifica su estado. Los
mecanismos de observación son completamente externos:

\begin{itemize}
    \item \textbf{Filesystem}: mediante \texttt{container.diff()} del Docker
    SDK, que retorna todos los archivos modificados, creados o eliminados
    respecto a la imagen original, sin montar volúmenes en el contenedor.

    \item \textbf{Procesos}: mediante \texttt{container.top()} con argumentos
    \texttt{-o pid,euid,ppid,cmd}, que lee la información de procesos desde el
    kernel del host sin ejecutar nada dentro del contenedor.

    \item \textbf{Logs de aplicación}: mediante lectura directa de archivos de
    log en el host, escritos por el contenedor a través de volúmenes montados en
    rutas explícitamente definidas en el YAML del escenario.

    \item \textbf{Red}: mediante un contenedor \textit{sidecar} de captura
    (basado en \texttt{nicolaka/netshoot}) que comparte el \textit{network
    namespace} del contenedor vulnerable. El \textit{sidecar} ejecuta
    \texttt{tcpdump} y vuelca el resultado a un archivo en un \textit{bind
    mount} del host. El monitor de red lee ese archivo desde fuera del
    contenedor vulnerable y emite los eventos correspondientes. El contenedor
    vulnerable nunca interactúa con \texttt{tcpdump} ni sabe que su tráfico se
    está observando.
\end{itemize}

Los cuatro monitores son coordinados por un \textit{orquestador} que provee un
único ciclo de \textit{polling} por entorno, agrupa los eventos emitidos y los
entrega en orden cronológico a la Capa 3. El detalle de esta coordinación,
incluyendo el orquestador global que supervisa varios entornos en paralelo,
corresponde al capítulo \ref{sec:implementacion} de implementación.

\subsubsection{Capa 3: Motor de Evaluación}

Consume los eventos generados por la Capa 2 y evalúa condiciones de victoria
definidas mediante reglas declarativas en archivos YAML. Determina cuándo un
participante ha completado exitosamente un ejercicio sin requerir la entrada
manual de \textit{flags}. Esta capa no tiene acceso directo a los contenedores:
todo su conocimiento del mundo exterior llega exclusivamente por los eventos de
Capa 2.

\subsubsection{Capa 4: Administración}

Gestiona el ciclo de vida de los entornos: crear, reiniciar o destruir según sea
necesario. Mantiene registro del progreso de todos los participantes según lo
indicado por la Capa 3 y proporciona interfaces de línea de comandos para
educadores y estudiantes.

\section{Contratos entre Capas}

Los contratos definen qué datos cruzan cada frontera, en qué formato, quién los
produce y quién los consume. El objetivo es que cada capa pueda desarrollarse,
testearse y reemplazarse de forma independiente siempre que respete estos
contratos.

\subsection{Capa 1 → Capa 2: Estado Observable}

La Capa 1 no comunica nada activamente a la Capa 2. Este es un contrato
implícito: la Capa 2 lee el estado de la Capa 1, y la Capa 1 no sabe que está
siendo observada. La interfaz entre ambas es el estado observable del entorno
virtualizado, accedido a través de tres canales descritos en la Sección
anterior.

\begin{sloppypar}
    Toda interacción con los contenedores se canaliza a través de la interfaz
    abstracta \texttt{EnvironmentProvider}, que encapsula la tecnología de
    virtualización:
\end{sloppypar}

\begin{verbatim}
class EnvironmentProvider(ABC):
    def create(self, scenario_id, participant_id, image,
               attacker_image=None, ports=None, attacker_ports=None,
               timeout_seconds=1800, startup_delay=None,
               enable_network_capture=False,
               capture_filter="") -> Environment
    def destroy(self, env_id) -> None
    def reset(self, env_id) -> None
    def get_status(self, env_id) -> Environment
    def exec_run(self, env_id, cmd, ...) -> ExecResult
    def get_filesystem_diff(self, env_id) -> list[dict]
    def get_processes(self, env_id) -> list[ProcessInfo]
    def cleanup_all(self) -> None
\end{verbatim}

En el estado actual del sistema la única implementación concreta de esta
interfaz es \texttt{DockerRootlessProvider}, que utiliza el SDK oficial de
Docker (\texttt{docker-py}) sobre un \textit{daemon} en modo \textit{rootless}.
Este proveedor materializa cada \texttt{Environment} como un conjunto de
recursos Docker etiquetados con \texttt{flexipwn.managed=true} y
\texttt{flexipwn.env\_id}. Mantiene los nombres de contenedores y redes
derivados del \texttt{env\_id} para permitir reconciliación tras una caída del
proceso. Una implementación alternativa (\texttt{PodmanProvider},
\texttt{KVMProvider}, etc.) solo debe respetar la firma de la interfaz para que
las capas superiores puedan operarla sin cambios.

\subsection{Capa 2 → Capa 3: Eventos Estructurados}

La Capa 2 produce eventos JSON estructurados que la Capa 3 consume. Este es el
contrato más crítico del sistema: define el lenguaje común entre la observación
y la evaluación. Todos los eventos comparten un envelope común validado por
Pydantic v2:

\begin{verbatim}
class MonitorEvent(BaseModel):
    timestamp: datetime
    monitor_type: Literal["filesystem", "process", "log", "network"]
    event_type: str
    env_id: str
    participant_id: str
    scenario_id: str
    details: dict[str, Any]
\end{verbatim}

El campo \texttt{monitor\_type} identifica la fuente del evento y
\texttt{event\_type} especifica la naturaleza del cambio. El campo
\texttt{details} varía según el monitor que produce el evento.

\subsubsection*{Eventos de filesystem}

Generados por el monitor de filesystem a partir de \texttt{container.diff()}.
Los \texttt{event\_type} posibles y los campos de \texttt{details} asociados
son:

\begin{itemize}
    \item \textbf{\texttt{file\_created}}: archivo nuevo detectado respecto a la
    imagen original. \texttt{details} contiene \texttt{path} (ruta absoluta) y
    \texttt{kind} (valor 1, según la convención de \texttt{container.diff()}).

    \item \textbf{\texttt{file\_modified}}: archivo existente en la imagen que
    fue modificado. \texttt{details} contiene \texttt{path} y \texttt{kind}
    (valor 0).

    \item \textbf{\texttt{file\_deleted}}: archivo de la imagen original
    eliminado durante el ejercicio. \texttt{details} contiene \texttt{path} y
    \texttt{kind} (valor 2). Cabe notar que este tipo de evento no tiene un
    target correspondiente en el schema YAML: ningún escenario del sistema
    evalúa la eliminación de archivos como condición de éxito en el estado
    actual.
\end{itemize}

\subsubsection*{Eventos de procesos}

Generados por el monitor de procesos a partir de \texttt{container.top()}. El
único \texttt{event\_type} implementado es:

\begin{itemize}
    \item \textbf{\texttt{process\_spawned}}: proceso nuevo detectado respecto
    al \textit{baseline} capturado al iniciar el monitor. Los campos de
    \texttt{details} son \texttt{pid}, \texttt{euid} (effective UID del
    proceso), \texttt{ppid}, \texttt{cmd} (comando completo), \texttt{lstart}
    (timestamp de inicio del proceso en formato \texttt{"Mon Oct 23 10:25:44
    2023"}), \texttt{ppid\_cmd} (comando del proceso padre directo),
    \texttt{ancestor\_cmds} (lista de comandos de los ancestros, del más cercano
    al más lejano) y \texttt{process\_id} (hash SHA-256 truncado a 12 caracteres
    de la combinación \texttt{pid:lstart}, que identifica claramente al proceso
    incluso ante reutilización de PID gracias al timestamp de inicio, pero
    cuando \texttt{lstart} no está disponible se recurre a \texttt{pid:cmd}).
    Los campos \texttt{ppid\_cmd} y \texttt{ancestor\_cmds} son los que
    alimentan los filtros \texttt{ppid\_cmd\_contains} y
    \texttt{ancestor\_contains} de la Capa 3.
\end{itemize}

\subsubsection*{Eventos de logs}

Generados por el monitor de logs mediante lectura directa de archivos en el host
escritos a través de volúmenes Docker.
\begin{sloppypar}
    \begin{itemize}
        \item \textbf{\texttt{log\_entry}}: línea de log nueva detectada en un
        archivo monitoreado. El monitor autodetecta el formato de cada línea, lo
        que produce dos variantes de \texttt{details}. Si la línea parsea como
        JSON, contiene \texttt{source\_file} (ruta del archivo de log en el
        host) y \texttt{parsed} (objeto JSON completo de la línea). Si no, la
        línea se trata como texto plano y \texttt{details} contiene
        \texttt{source\_file} y \texttt{raw\_line} (la línea íntegra). Esta
        segunda variante es la que consumen los escenarios basados en MySQL,
        cuyo \textit{query log} es texto plano.
    \end{itemize}
\end{sloppypar}

\subsubsection*{Eventos de red}

Generados por el \texttt{NetworkMonitor} a partir de la captura de tráfico
realizada por el contenedor \textit{sidecar} basado en
\texttt{nicolaka/netshoot} (descrito en la Sección~\ref{sec:capa2}). Se emiten
dos tipos de eventos:

\begin{itemize}
    \item \textbf{\texttt{network\_payload}}: paquete observado con contenido
    ASCII imprimible relevante. \texttt{details} contiene \texttt{data} (texto
    ASCII imprimible extraído del payload, la captura se realiza a nivel de
    paquete con \texttt{tcpdump -A}, y el parser observa además las banderas TCP
    del paquete, por ejemplo \texttt{[S.]} o \texttt{[P.]}, para clasificarlo),
    \texttt{src\_ip}, \texttt{src\_port}, \texttt{dst\_ip} y \texttt{dst\_port}.

    \item \textbf{\texttt{network\_connection}}: conexión TCP establecida hacia
    un destino, detectada por el primer paquete \texttt{SYN-ACK} (\texttt{Flags
    [S.]}). \texttt{details} contiene \texttt{src\_ip}, \texttt{dst\_ip} y
    \texttt{dst\_port}.
\end{itemize}

La Capa 2 garantiza que todo evento emitido cumple el schema antes de enviarse,
que los eventos se emiten en orden cronológico por monitor, y que no se emiten
duplicados para el mismo cambio de estado.

\subsection{Capa 3 → Capa 4: Resultados de Evaluación}

La Capa 3 comunica a la Capa 4 los resultados de la evaluación de condiciones de
éxito mediante dos estructuras Pydantic complementarias.

\texttt{TargetResult} representa el estado de evaluación de un objetivo
individual:

\begin{verbatim}
class TargetResult(BaseModel):
    target_index: int
    target_type: str
    description: str
    matched: bool = False
    matched_at: datetime | None = None
    trigger_event: MonitorEvent | None = None
    children: list["TargetResult"] | None = None
\end{verbatim}

El campo \texttt{trigger\_event} contiene el \texttt{MonitorEvent} exacto que
causó el \textit{match}, lo que permite al educador reconstruir no solo qué
objetivo se completó, sino qué acción concreta del estudiante lo disparó. Esta
trazabilidad es uno de los valores pedagógicos diferenciales del sistema
respecto a los sistemas de \textit{flags}.

El campo \texttt{children} se utiliza únicamente cuando el target es un nodo
lógico (\texttt{and}, \texttt{or} o \texttt{not}). Contiene los
\texttt{TargetResult} de los sub-targets evaluados recursivamente, lo que
permite a la Capa 4 mostrar al educador no solo qué rama lógica matcheó, sino
qué sub-condiciones internas fueron las que la satisficieron.

\texttt{EvaluationResult} agrega el estado de todos los targets de un escenario
para un participante:

\begin{verbatim}
class EvaluationResult(BaseModel):
    scenario_id: str
    participant_id: str
    env_id: str
    condition: str                        # "any" | "all"
    targets: list[TargetResult]
    completed: bool
    completed_at: datetime | None = None
    progress: float                       # hojas matcheadas / total hojas
\end{verbatim}

El campo \texttt{progress} se calcula como la razón entre el número de
\emph{hojas} matcheadas y el total de hojas del árbol de targets del escenario.
Las hojas son los targets atómicos (\texttt{file\_created},
\texttt{process\_running}, \texttt{network\_payload}, etc.); los nodos lógicos
(\texttt{and}, \texttt{or}, \texttt{not}) no cuentan como hoja: aportan
estructura pero no se incluyen en el denominador del progreso. El campo
\texttt{completed} es \texttt{True} si y solo si la condición (\texttt{any} o
\texttt{all}) se satisface con el conjunto de targets de primer nivel marcados
como \texttt{matched}. Cuando un target de primer nivel es un nodo lógico, su
\texttt{matched} se determina recursivamente: \texttt{and} matchea si todos sus
hijos matchean, \texttt{or} si al menos uno, \texttt{not} si su único hijo no
matchea.

El \texttt{EvaluationEngine} invoca el callback \texttt{on\_update} cada vez que
un target cambia de estado o cuando \texttt{completed} pasa de \texttt{False} a
\texttt{True}. Los logros son irreversibles dentro de un \textit{run}: una vez
que \texttt{matched} es \texttt{True}, no vuelve a \texttt{False}. La Capa 4
consume estos resultados para persistirlos en SQLite mediante la entidad
\texttt{TargetResult} y para emitirlos a través de un \texttt{NotificationSink}
central que decide, según política por tipo, qué se imprime en el TTY del
educador y qué queda únicamente en log y en las vistas agregadas \texttt{dashboard} 
y \texttt{feed}. A partir de los instantes en que cada hoja
matchea, la Capa 4 deriva un registro de tiempos por run: la duración total y el
tiempo de cada etapa hasta su objetivo, que se ofrece al educador en esas mismas
vistas sin necesidad de almacenar datos adicionales.
\subsection{Capa 4 → Capa 1: Ciclo de Vida}

La Capa 4 es la única capa que modifica el estado de la Capa 1, y lo hace
exclusivamente a través de los métodos de la interfaz
\texttt{EnvironmentProvider}: la materialización de un entorno, su reinicio, su
detención o destrucción, y la reconciliación del estado persistido con los
recursos reales del proveedor. Estas operaciones son invocaciones internas que
la Capa 4 encadena para sostener los flujos expuestos al usuario; el detalle de
cómo cada acción del educador o del estudiante se traduce en una secuencia de
invocaciones al \texttt{EnvironmentProvider} se describe en el capítulo de
implementación.

Las consultas de estado (\texttt{get\_status}, \texttt{get\_processes},
\texttt{get\_filesystem\_diff}) son de solo lectura y, aunque la Capa 4 también
las utiliza para producir vistas al usuario, conviven con el uso que de ellas
hace la Capa 2 para generar eventos. La separación se preserva porque ambas
capas consumen información sin modificar el entorno.

\subsection{Educador → Capa 3: Archivo YAML del Escenario}

El educador define un escenario mediante un único archivo YAML que la Capa 3
parsea y valida con Pydantic v2 antes de cualquier ejecución. Este archivo es la
única pieza de configuración que el educador necesita escribir, además de la
imagen Docker del sistema vulnerable.

La estructura completa del archivo es la siguiente:

\begin{verbatim}
title: str               # nombre del ejercicio (requerido)
description: str         # descripción del ejercicio (requerido)
author: str              # autor del escenario (requerido)
level: str               # beginner | intermediate | advanced
category: str            # pwning | web | database | forensics | reversing

environment:
  image: str             # imagen Docker para Capa 1a (requerido)
  attacker_image: str    # imagen Docker para Capa 1b (opcional)
  log_paths: list[str]   # rutas de logs a monitorear (opcional)
  ports: list[str]       # puertos extra del vulnerable
                         # "host:container" (opcional)
  capture_filter: str    # filtro para tcpdump, ej. "port 3306"
                         # (opcional, solo aplica si hay targets
                         # de tipo network_*)
  startup_delay_seconds: float  # delay opcional antes de capturar
                         # el baseline si no hay HEALTHCHECK
                         # (null usa el default global)

hints:
  - str                  # pistas ordenadas para el estudiante (opcional)

targets:                 # lista de objetivos (requerido, mínimo 1)
  - type: str
    description: str
    # campos específicos según el tipo de target

condition: str           # any | all (requerido, aplica solo a la
                         # lista plana de targets, se ignora cuando
                         # los targets de primer nivel son nodos
                         # lógicos and/or/not)
timeout_seconds: int     # tiempo máximo en segundos (default: 1800)
\end{verbatim}

Los tipos de target disponibles se dividen en \emph{tipos atómicos}, que evalúan
una condición concreta sobre eventos de la Capa 2, y \emph{nodos lógicos}, que
componen otros targets en árboles recursivos.

\paragraph{Tipos atómicos}\mbox{}\\[0.0em]
\begin{description}
    \item[\texttt{file\_created}] Detecta la creación de un archivo nuevo.
    Requiere \texttt{path} (ruta absoluta o directorio terminado en \texttt{/}).
    Acepta \texttt{pattern} (glob) cuando \texttt{path} es un directorio.

    \item[\texttt{file\_modified}] Detecta la modificación de un archivo
    existente en la imagen original. Requiere \texttt{path} (ruta absoluta).

    \item[\texttt{process\_running}] Detecta un proceso activo con
    características específicas. Requiere \texttt{euid} (effective UID esperado;
    0 para root) y \texttt{cmd\_contains} (substring en el comando del proceso).
    Acepta \texttt{ppid\_cmd\_contains} (substring en el comando del proceso
    padre) y \texttt{ancestor\_contains} (substring en el comando de cualquier
    ancestro).

    \item[\texttt{log\_pattern}] Detecta una entrada de log que cumpla
    condiciones sobre campos parseados. Requiere \texttt{field\_matches}
    (diccionario de campo → regex). El campo especial \texttt{raw\_line} permite
    matchear contra la línea completa cuando el log es texto plano. Requiere que
    \texttt{environment.log\_paths} declare al menos una ruta.

    \item[\texttt{network\_payload}] Detecta payloads de paquetes de red cuyo
    contenido ASCII imprimible coincida con una expresión regular. Requiere
    \texttt{field\_matches} con la clave \texttt{data}, y el matching se realiza
    con \texttt{re.IGNORECASE} sobre el texto extraído del paquete. El
    \textit{sidecar} de captura comparte el \textit{network namespace} del
    contenedor vulnerable. Estos escenarios exigen definir un entorno atacante
    (\texttt{attacker\_image}): es el espacio controlado desde el cual el
    estudiante lanza el ataque y, por tanto, la fuente del tráfico que el
    monitor observa. Opcionalmente se acota la captura con
    \texttt{capture\_filter} en el bloque \texttt{environment}.

    \item[\texttt{network\_connection}] Detecta conexiones TCP establecidas
    (paquete \texttt{SYN-ACK}) hacia un destino. Requiere \texttt{dst\_port};
    acepta \texttt{dst\_ip} para restringir a una IP específica.
\end{description}

\paragraph{Nodos lógicos}\mbox{}\\[1.0em]
Los nodos lógicos componen otros targets en árboles recursivos y se identifican
por su \texttt{type}. Sus sub-targets se definen en el campo \texttt{targets},
que puede contener tanto tipos atómicos como otros nodos lógicos.

\begin{description}
    \item[\texttt{and}] Matchea cuando todos sus sub-targets matchean. Requiere
    al menos dos sub-targets.

    \item[\texttt{or}] Matchea cuando al menos un sub-target matchea. Requiere
    al menos dos sub-targets.

    \item[\texttt{not}] Matchea cuando su único sub-target no matchea. Requiere
    exactamente un sub-target. No puede aparecer como target de primer nivel del
    escenario, ya que un escenario que se completara por la ausencia de un
    evento podría darse por exitoso sin acción alguna del estudiante; este caso
    se valida al cargar el YAML.
\end{description}

El campo de primer nivel \texttt{condition} (\texttt{any} o \texttt{all}) actúa
como azúcar sintáctica para listas planas de targets atómicos: es equivalente a
envolver la lista en un nodo \texttt{or} o \texttt{and} respectivamente.

A modo de ejemplo, el siguiente archivo YAML define un escenario de escalación
de privilegios mediante \texttt{sudo vim}:

\begin{verbatim}
title: "Privilege escalation via sudo vim"
description: >
  El usuario ctfuser tiene permisos sudo sobre vim sin contraseña.
  Escala privilegios y crea evidencia de acceso root.
author: "Dylan Riquelme"
level: beginner
category: pwning

environment:
  image: "flexipwn/vuln-sudo"
  attacker_image: "flexipwn/attacker"
  log_paths: []
  ports: []
  startup_delay_seconds: 3

hints:
  - "Conéctate al vulnerable desde el atacante: ssh ctfuser@flexipwn-<env_id>-vulnerable (clave: ctfpassword)"
  - "Una vez dentro, revisa tus permisos sudo con: sudo -l"
  - "Vim puede ejecutar comandos del sistema con :!comando"

targets:
  - type: file_created
    path: "/root/"
    pattern: "*.txt"
    description: "Archivo .txt creado en directorio root"

  - type: file_modified
    path: "/etc/passwd"
    description: "Archivo de usuarios del sistema modificado"

  - type: process_running
    euid: 0
    cmd_contains: "bash"
    ppid_cmd_contains: "vim"
    ancestor_contains: "sudo"
    description: "Shell root lanzada desde vim (escalación exitosa)"

condition: all
timeout_seconds: 1800
\end{verbatim}

El schema se valida al momento de cargar el escenario, antes de cualquier
ejecución. Los errores de validación producen mensajes que indican el campo
problemático y el valor esperado, sin requerir ejecutar el sistema para
descubrirlos.

Más allá de la carga de escenarios, este contrato también da soporte a las
operaciones de validación y limpieza que el educador necesita sobre los
escenarios y la actividad acumulada. Se exponen más detalles en el capítulo de
implementación.

% TODO figura: diagrama de flujo de datos entre las cuatro capas + Educador
% (estado observable, eventos estructurados, resultados de evaluación, ciclo de
% vida y YAML del escenario) (ver Paso 4)

\section{Modelo de Datos}

El estado del sistema se persiste en una base de datos SQLite gestionada
mediante SQLModel. La elección de SQLite responde al contexto de despliegue: una
instalación autohospedada de escala pequeña a mediana que no justifica la
operación de un servidor de base de datos independiente.

El modelo contempla cinco entidades principales. Todas usan identificadores
\texttt{UUID} como clave primaria, generados automáticamente al persistir, lo
que evita filtrar información temporal en los identificadores y simplifica la
fusión de bases de datos entre instalaciones distintas.

\textbf{Scenario} almacena los metadatos del escenario cargado por el educador,
incluyendo la ruta al archivo YAML, el contenido íntegro del YAML para
reproducibilidad, y campos desnormalizados (\texttt{title},
\texttt{description}, \texttt{level}, \texttt{category}, \texttt{image},
\texttt{attacker\_image}, \texttt{timeout\_seconds}) para consultas directas sin
reparsear el YAML.

\textbf{Participant} representa a un estudiante del sistema. Tiene
\texttt{username} con índice único, evitando colisiones a nivel de base de
datos. La contraseña se persiste como \textit{hash} bcrypt, el texto plano solo
se expone una vez al momento de crear o resetear el participante, y nunca se
vuelve a recuperar.

\textbf{ExerciseRun} vincula a un participante con un escenario y registra el
estado del ejercicio. Mantiene índices ordinarios sobre las claves foráneas
\texttt{scenario\_id} y \texttt{participant\_id} para acelerar las consultas
frecuentes (runs por participante, runs por escenario), y un índice único sobre
\texttt{env\_id} para garantizar que el identificador de entorno expuesto al
participante no se repita entre runs. Persiste además los datos efímeros de
acceso del estudiante (\texttt{attacker\_ssh\_port},
\texttt{attacker\_ssh\_username}, \texttt{attacker\_ssh\_password}) para que el
educador pueda recuperar las credenciales en cualquier momento desde la CLI. La
contraseña SSH se guarda en texto plano dentro del archivo SQLite del educador,
con la trazabilidad documentada en el código sobre por qué este compromiso es
aceptable para el contexto educativo. El campo \texttt{status} recorre los
estados de \texttt{RunStatus} a lo largo del ciclo de vida del ejercicio (desde
\texttt{pending} hasta los estados terminales \texttt{completed},
\texttt{failed}, \texttt{timeout} o \texttt{stopped}, pasando por
\texttt{running}, \texttt{stopping} y \texttt{resetting}). Mientras que los
campos \texttt{reset\_payload} y \texttt{daemon\_message} acompañan a esas
transiciones mediadas por el daemon. La validación de unicidad sobre el par
(\texttt{scenario\_id}, \texttt{participant\_id}) para runs activos se hace a
nivel aplicación, no de base de datos: el mismo participante puede tener
múltiples runs históricos del mismo escenario, lo que permite reintentos sin
borrar el historial previo.

\textbf{TargetResult} registra el estado de evaluación de cada objetivo
individual dentro de un \textit{run}, incluyendo el evento exacto que causó el
\textit{match} serializado como JSON. Cada fila corresponde a una \emph{hoja}
del árbol de targets: el árbol del escenario se aplana en orden DFS y cada hoja
recibe un \texttt{target\_index} secuencial. Los nodos lógicos (\texttt{and},
\texttt{or}, \texttt{not}) no se persisten como filas, aportan estructura solo
en memoria. El reenlace entre la hoja del modelo Pydantic de la Capa 3 y su fila
es posicional por \texttt{target\_index} en ese mismo orden DFS, lo que evita
guardar un camino jerárquico en la base de datos. Sus columnas son \texttt{id},
\texttt{run\_id}, \texttt{target\_index}, \texttt{target\_type},
\texttt{description}, \texttt{matched}, \texttt{matched\_at},
\texttt{trigger\_event} y \texttt{reset\_at}. Este último permite reintentos sin
borrar el historial, ya que el intento actual es el conjunto de filas con
\texttt{reset\_at} nulo. Tiene índice ordinario sobre \texttt{run\_id} para
recuperar eficientemente todos los resultados de un run.

\textbf{RunEvent} guarda el flujo completo de eventos del monitor asociados a un
run (\texttt{monitor\_type}, \texttt{event\_type}, \texttt{details\_json},
\texttt{timestamp}). Sirve como bitácora auditable del ejercicio, independiente
de qué targets matcheen o no: permite reconstruir a posteriori la secuencia de
acciones del estudiante. Tiene índice ordinario sobre \texttt{run\_id}.

% TODO: copiar flexipwn_erd_v2.pdf desde Notion a imagenes/erd.pdf
% y descomentar el bloque siguiente para incluir el diagrama.
% \begin{figure}[h]
% \centering
% \includegraphics[width=0.85\textwidth]{imagenes/erd.pdf}
% \caption{Diagrama entidad-relación del modelo de datos de FlexiPwn.}
% \label{fig:erd}
% \end{figure}

\section{Aislamiento de Entornos}

El aislamiento entre los entornos de práctica de distintos participantes opera
en dos dimensiones complementarias: aislamiento de red y aislamiento de
filesystem.

\subsection*{Modelo de red de doble capa}

Cada entorno instancia dos redes Docker independientes con roles diferenciados.
La red interna (\texttt{flexipwn-\{env\_id\}}), creada con
\texttt{internal=True}, conecta el sistema vulnerable y el sistema atacante
entre sí sin permitir tráfico hacia el exterior. Esta configuración garantiza
que un ejercicio de explotación activa no genere tráfico hacia redes externas ni
hacia otros entornos del sistema. La red externa
(\texttt{flexipwn-\{env\_id\}-ext}), un bridge estándar sin restricciones,
conecta exclusivamente el contenedor atacante al host. Esto permite que el
participante acceda al entorno atacante mediante SSH desde su máquina,
manteniendo al mismo tiempo el contenedor vulnerable completamente aislado del
exterior.

El contenedor atacante opera en ambas redes simultáneamente: su interfaz en la
red interna le permite resolver el nombre DNS del vulnerable
(\texttt{flexipwn-\{env\_id\}-vulnerable}) y establecer conexiones hacia él. Su
interfaz en la red externa permite que los port bindings definidos en el YAML
del escenario sean alcanzables desde el host. El contenedor vulnerable solo
tiene presencia en la red interna.

El contenedor sidecar de captura de red (Sección~\ref{sec:capa2}) no se conecta
a ninguna red Docker directamente: comparte el namespace de red del contenedor
vulnerable mediante \texttt{--network container:\{nombre\_vulnerable\}}, lo que
le permite observar todo el tráfico del vulnerable sin introducir interfaces
adicionales en la topología de red del entorno.

% TODO figura: topología de red (red interna/externa + sidecar, 3 contenedores)
% — vulnerable solo en la red interna, atacante en ambas redes, sidecar
% compartiendo el netns del vulnerable (ver Paso 4)

\subsection*{Aislamiento de filesystem}

Los bind mounts efímeros descritos en la Sección~\ref{sec:implementacion} se
crean en su mayoría con permisos \texttt{0o700}, accesibles únicamente por el
proceso del daemon rootless. La base de volúmenes y el directorio de captura de
red usan ese modo. Una excepción son los directorios de logs que se montan al
contenedor vulnerable, creados con \texttt{0o777} para que el proceso del
contenedor pueda escribir sus archivos de log. Al destruir el entorno, el árbol
de directorios completo se elimina mediante \texttt{shutil.rmtree()}. El
contenedor \textit{sidecar} de captura se detiene y elimina al inicio del ciclo
de destrucción, antes de detener el contenedor vulnerable y atacante, y su
directorio de captura se borra en ese mismo paso.

\chapter{Implementación}

Este capítulo describe las decisiones técnicas tomadas durante la
implementación de FlexiPwn, organizadas según las capas de la arquitectura.

\section{Stack Tecnológico}

% [Pegar aquí el texto del documento stack-tecnologico-flexipwn.pdf
% que ya tienes generado. Incluye: lenguaje Python 3.12, uv,
% decisión por cada capa, tabla resumen al final]

\section{Decisión de Tecnología de Virtualización}

% [Pegar aquí el texto del documento virt-decision-flexipwn.pdf,
% incluyendo la tabla comparativa y la justificación de Docker rootless]

\section{Capa 1: Orquestador de Contenedores}

% [Descripción de DockerRootlessProvider, detección de socket,
% estrategia de baseline con healthcheck, rollback en create()]

\section{Capa 2: Monitoreo Pasivo}

% [Descripción de FilesystemMonitor con container.diff(),
% ProcessMonitor con container.top() y lstart,
% MonitorOrchestrator, y baseline de startup]

\subsection{Limitaciones del Monitoreo por Polling}

El monitor de procesos consulta \texttt{container.top()} cada dos segundos.
Procesos que nacen y terminan dentro de ese intervalo son invisibles para
el sistema. Los vectores no detectables en el estado actual incluyen:
procesos efímeros con duración menor a dos segundos, técnicas de
\textit{PPID spoofing}, y ataques \textit{fileless} mediante llamadas al
sistema como \texttt{ptrace} o \texttt{memfd\_create}. Para el alcance
educativo del proyecto, este nivel de detección es aceptable dado que los
ejercicios no requieren detectar procesos sub-segundo ni técnicas de evasión
avanzadas esperables en un atacante universitario.

De forma análoga, \texttt{container.diff()} retorna una instantánea del
estado acumulado del filesystem desde el inicio del contenedor. Para
distinguir cambios del sistema de inicialización de cambios realizados por
el estudiante, se captura un \textit{baseline} del \textit{diff} una vez
que el contenedor está estable. Si la imagen tiene un \texttt{HEALTHCHECK}
configurado, el sistema espera hasta que el contenedor reporte estado
\texttt{healthy} antes de capturar el baseline; de lo contrario, espera un
delay configurable (por defecto 3 segundos) como mecanismo de respaldo.

\subsection{Monitor de Red}
\label{sec:capa2}

Para escenarios que requieren observar tráfico de red, la Capa 2 instancia 
un cuarto monitor: \texttt{NetworkMonitor}. A diferencia de los tres monitores 
anteriores, este monitor depende de un contenedor sidecar que comparte el 
namespace de red del sistema vulnerable.

\subsubsection*{Contenedor sidecar}
\begin{sloppypar}
    Al crear el entorno, si el escenario define al menos un target de tipo 
    \texttt{network\_*}, \texttt{DockerRootlessProvider} lanza un contenedor 
    adicional basado en la imagen \texttt{nicolaka/netshoot} con la siguiente 
    configuración:
\end{sloppypar}

\begin{itemize}
    \item \textbf{Namespace de red compartido}: mediante 
    \texttt{--network container:\{vulnerable\}}, el sidecar comparte 
    exactamente el mismo namespace de red que el vulnerable. Esto hace 
    visible todo el tráfico del vulnerable desde dentro del sidecar, 
    incluyendo comunicaciones por la interfaz de loopback 
    (\texttt{127.0.0.1}), que de otro modo no serían capturables desde 
    una perspectiva externa.
    
    \item \textbf{Comando de captura}: \texttt{tcpdump -i any -A -n -l}, 
    con un filtro opcional definido en el campo 
    \texttt{environment.capture\_filter} del YAML del escenario 
    (por ejemplo, \texttt{port 3306} para restringir la captura 
    a tráfico MySQL).
    
    \item \textbf{Salida en bind mount}: la captura se escribe en un 
    archivo \texttt{traffic.txt} en un directorio del host con 
    permisos \texttt{0o700}, al que solo tiene acceso el proceso de 
    FlexiPwn.
\end{itemize}

El uso de \texttt{-n} desactiva la resolución de nombres, forzando IPs 
numéricas en el output y evitando que el parser del monitor falle ante 
hostnames como \texttt{localhost}. El flag \texttt{-l} activa el 
line-buffering de tcpdump, reduciendo la latencia entre la escritura 
del paquete y su disponibilidad en el archivo del host.

\subsubsection*{Parsing packet-based}

\texttt{NetworkMonitor} lee el archivo \texttt{traffic.txt} de forma 
incremental en cada ciclo de polling, procesando el output de tcpdump 
como unidades de paquete completo en lugar de línea por línea. Esta 
decisión es necesaria porque el output de \texttt{tcpdump -A} en modo 
\texttt{-i any} mezcla bytes ASCII del payload con bytes binarios de 
los headers TCP/IP, y el carácter \texttt{0x0a} (LF) puede aparecer 
dentro de los headers, partiendo artificialmente el payload en múltiples 
líneas si se procesa con \texttt{splitlines()}.

El monitor detecta las cabeceras de paquete mediante una expresión regular 
que acepta el formato de \texttt{-i any} (que incluye el nombre de interfaz 
y dirección entre el timestamp y la IP) y extrae de ellas las IPs y 
puertos de origen y destino. Las líneas entre cabeceras se acumulan en un 
buffer de payload. Cuando se detecta una nueva cabecera, el buffer del 
paquete anterior se procesa: se extrae todo el texto en rango ASCII 
imprimible (\texttt{0x20--0x7e}), y si el resultado supera los cuatro 
caracteres se emite un \texttt{MonitorEvent} de tipo \texttt{network\_payload} 
con el texto extraído y los metadatos del paquete. Adicionalmente, si la 
cabecera contiene \texttt{Flags [S.]}, se emite un evento 
\texttt{network\_connection} con la IP y puerto de destino.

Una implicación importante del modelo packet-based es que el último 
paquete de una ráfaga de tráfico corta podría quedar atrapado en el 
buffer si no llega una cabecera posterior que dispare su procesamiento. 
Para evitar esto, al final de cada ciclo de polling el monitor procesa 
el paquete en curso si el buffer no está vacío, sin esperar al siguiente 
ciclo.

\section{Capa 3: Motor de Evaluación}

% [Descripción del EvaluationEngine, tipos de target implementados,
% condiciones any/all, evaluadores de Capa 3]

Los tipos de target implementados en el estado actual del sistema son:

\begin{description}
    \item[\texttt{file\_created}, \texttt{file\_modified}, \texttt{file\_exists}] 
    Operan sobre eventos del monitor de filesystem. Permiten detectar 
    creación o modificación de archivos en rutas absolutas o mediante 
    patrones glob.
    
    \item[\texttt{process\_running}] Opera sobre eventos del monitor de 
    procesos. Detecta procesos activos por effective UID y substring 
    en el comando.
    
    \item[\texttt{log\_pattern}] Opera sobre eventos del monitor de logs. 
    Aplica un diccionario de campo→regex sobre los campos de la entrada 
    de log parseada como JSON.
    
    \item[\texttt{network\_payload}] Opera sobre eventos \texttt{network\_payload} 
    del \texttt{NetworkMonitor}. Aplica un diccionario de campo→regex 
    sobre el campo \texttt{data} del evento, que contiene el texto 
    ASCII imprimible extraído del payload del paquete. El matching se 
    realiza con \texttt{re.IGNORECASE}, de modo que el educador no 
    necesita contemplar variaciones de capitalización en los patrones.
    
    \item[\texttt{network\_connection}] Opera sobre eventos 
    \texttt{network\_connection} del \texttt{NetworkMonitor}. Detecta 
    conexiones TCP establecidas (paquete SYN-ACK) hacia un puerto de 
    destino específico, con filtro opcional por IP de destino.
    
    \item[\texttt{and}, \texttt{or}, \texttt{not}] Nodos lógicos 
    recursivos que permiten componer condiciones complejas sobre 
    cualquier combinación de targets atómicos. El campo 
    \texttt{condition: any|all} del escenario actúa como azúcar 
    sintáctica para listas planas de targets sin nodos lógicos.
\end{description}

Los tipos \texttt{http\_response\_contains} y 
\texttt{database\_query\_result} están definidos en el schema Pydantic 
pero no tienen evaluador implementado. Quedan documentados como trabajo 
futuro; el segundo en particular fue reemplazado funcionalmente por 
\texttt{network\_payload} con filtro \texttt{capture\_filter: "port 3306"}, 
que captura las queries directamente del tráfico TCP sin requerir acceso 
a la base de datos.

\section{Capa 4: Administración}

% [Descripción de la CLI con Typer + Rich, comando demo privesc,
% MonitorOrchestrator como punto de integración]
\chapter{Validación}
\section{Evaluación}

% ===========================================================================
% BOSQUEJO DE TRABAJO — Por desarrollar a medida que se realicen pruebas
% reales del sistema (CLCERT, red DCC, escenarios multi-estudiante, etc.).
%
% Las secciones 5.1 a 5.3 se podrían añadir ya: los tres escenarios están
% funcionando en el repo y los walkthroughs ya están redactados en el README
% del repositorio público.
%
% Las secciones 5.4 y 5.5 dependen de pruebas que aún faltan.
% ===========================================================================

% ---------------------------------------------------------------------------
% 5.1 Metodología y criterios de validación
% ---------------------------------------------------------------------------
% - Recordar el criterio declarado en el Capítulo 1: validación automática
%   para al menos tres categorías entre web / database / pwning / forensics /
%   reversing, sin necesidad de flags manuales.
% - Describir cómo se valida cada escenario: setup (imagen y configuración),
%   ejecución del exploit por parte del estudiante, targets esperados,
%   targets disparados, evidencia capturada (panel SSH, output del REPL).
% - Aclarar que la validación opera en dos niveles:
%     (a) los escenarios funcionan y disparan los targets de forma correcta;
%     (b) el sistema se sostiene bajo uso realista — pruebas pendientes.

% ---------------------------------------------------------------------------
% 5.2 Cobertura del criterio de éxito
% ---------------------------------------------------------------------------
% - Tabla resumen con los tres escenarios:
%     columna 1: escenario
%     columna 2: categoría (privesc, sqli, network)
%     columna 3: monitores involucrados (FS, Proc, Log, Red)
%     columna 4: tipos de target ejercitados
% - Mostrar que entre los tres se cubren los cuatro monitores y los nodos
%   lógicos (al menos un escenario con and / or si aplica).
% - Conectar con el objetivo declarado: "al menos tres categorías".

% ---------------------------------------------------------------------------
% 5.3 Escenarios validados
% ---------------------------------------------------------------------------

% \subsection{Privilege Escalation vía sudo vim (pwning)}
%   - Setup: imagen flexipwn/vuln-sudo (sshd + sudo NOPASSWD vim).
%   - Walkthrough del estudiante:
%       ssh student-XXXXXX@host -p PORT       (al atacante)
%       ssh ctfuser@flexipwn-<env>-vulnerable  (al vulnerable)
%       sudo -l                                (descubre el privilegio)
%       sudo vim -c ':!bash'                   (escalada)
%       echo "pwned" > /root/flag.txt          (evidencia)
%       echo "..." >> /etc/passwd              (modificación)
%   - Targets disparados:
%       file_created /root/*.txt
%       file_modified /etc/passwd
%       process_running euid=0 cmd=bash ancestor_contains=sudo
%   - Salida esperada del REPL del educador.

% \subsection{SQL Injection Login Bypass (web)}
%   - Setup: imagen flexipwn/vuln-sqli-mysql (Flask + MySQL en :5001, general_log on).
%   - Walkthrough del estudiante:
%       payload 1: admin' OR '1'='1' --        (bypass clásico)
%       payload 2: UNION SELECT ... FROM sensitive_data --  (extracción)
%   - Targets disparados (mezcla log_pattern + network_payload con
%     capture_filter "port 3306"):
%       log_pattern OR.*1.*=.*1
%       log_pattern authentication_success
%       log_pattern SELECT.*sensitive_data
%       network_payload OR.*1.*=.*1
%       network_payload FLAG{sql_injection_detected}
%   - Discutir cómo logs y red se refuerzan: el mismo payload activa varios
%     targets por canales independientes.

% \subsection{Command Injection -> Reverse Shell (network)}
%   - Setup: imagen flexipwn/vuln-command-injection (Flask /ping?host= + nc).
%   - Walkthrough del estudiante (necesita dos sesiones SSH al atacante):
%       Terminal A: nc -lvp 4444
%       Terminal B: curl ".../ping?host=;nc <attacker> 4444 -e /bin/bash"
%       Terminal A (ya con la reverse shell): echo "owned" > /root/win.txt
%   - Targets disparados:
%       network_connection dst_port=4444
%       file_created /root/*.txt
%   - Nota: este escenario justifica el sidecar tcpdump (Capa 2) y la red
%     externa atacante-host (Capa 1).

% ---------------------------------------------------------------------------
% 5.4 Pruebas de robustez del sistema  — PENDIENTES
% ---------------------------------------------------------------------------
% Cada uno requiere ejecutar la prueba y registrar el resultado:
%
% - Manejo de YAML malformado: validar que Pydantic devuelve un mensaje claro
%   identificando el campo problemático sin ejecutar el sistema.
% - Manejo de contenedor caído a mitad del run: verificar que on_stopped lleva
%   al estado failed o stopped correctamente, y que el daemon no queda colgado.
% - Manejo de participante duplicado / username repetido.
% - Manejo de scenario id no encontrado en run start.
% - Pruebas de batch-start con N estudiantes simultáneos (aula completa):
%   contar timeouts, fallos por límite de puertos del rango 2200-2299,
%   contención de la base SQLite.
% - Tiempo de detección desde el evento (acción del estudiante) hasta la
%   alerta visible en el REPL del educador. Medir gap real de polling.
% - "Trampa" en tcpdump: enviar un payload SQLi por fuera del flujo normal del
%   escenario (por ejemplo, desde una terminal del host atacante con curl
%   directo) y verificar que el monitor de red lo detecta igualmente.
% - run reset: verificar que el env_id cambia, que el ExerciseRun se preserva,
%   que los TargetResult del intento previo quedan con reset_at, y que el
%   nuevo intento empieza limpio.
% - exit/Ctrl+D en cliente attached vs daemon stop: comprobar la separación.

% ---------------------------------------------------------------------------
% 5.5 Pruebas en infraestructura real  — PENDIENTES
% ---------------------------------------------------------------------------
% - Despliegue en servidor del CLCERT.
% - Exposición controlada a la red del DCC / la universidad.
% - Pruebas de aislamiento entre estudiantes:
%     un estudiante intenta resolver hostnames del entorno de otro
%     un estudiante intenta egresar al host
%     un estudiante intenta llegar al daemon FlexiPwn
% - Pruebas de escape de contenedor (al menos un par de técnicas conocidas):
%   verificar que el modo rootless contiene el escape como se afirmó en
%   el Capítulo de Implementación.
% - Sesión piloto con un grupo de estudiantes reales (si se llega a alcanzar):
%   medir UX, tiempos de onboarding, fricciones del flujo educador-estudiante.

% ---------------------------------------------------------------------------
% 5.6 Discusión  (cierre del capítulo)
% ---------------------------------------------------------------------------
% - Qué demuestran los tres escenarios respecto al criterio del Capítulo 1.
% - Qué queda como pendiente y cuál es su impacto.
% - Hilo con el Capítulo de Conclusión: trabajo futuro y reflexión final.

\chapter{Conclusión}

% ===========================================================================
% BOSQUEJO DE TRABAJO — Por desarrollar tras la validación del Capítulo 5.
%
% La mayoría de las secciones puede escribirse a partir de lo ya documentado
% (objetivos del Capítulo 1, contribuciones implícitas en Capítulos 3 y 4,
% limitaciones reconocidas en Capítulos 3 y 4). La sección 6.5 es opcional
% y depende de la reflexión personal una vez cerrado el resto.
% ===========================================================================

% ---------------------------------------------------------------------------
% 6.1 Síntesis del trabajo realizado
% ---------------------------------------------------------------------------
% - Resumir qué se construyó: plataforma educativa autohospedable de
%   ciberseguridad ofensiva con validación automática basada en estado del
%   sistema, motor de evaluación con targets atómicos y nodos lógicos
%   (and/or/not), y daemon REPL multi-cliente para uso educativo.
% - Recorrer los OBJETIVOS ESPECÍFICOS del Capítulo 1 y marcar el estado:
%     1. Arquitectura modular multi-entorno con virtualización   (cumplido).
%     2. Configuración declarativa YAML accesible al educador    (cumplido).
%     3. Motor de validación automática basado en estado         (cumplido).
%     4. Herramientas CLI de seguimiento y administración        (cumplido).
%     5. Poner la plataforma a disposición del CLCERT            (parcial:
%        depende de la entrega y las pruebas en su infraestructura — ver
%        Capítulo 5).

% ---------------------------------------------------------------------------
% 6.2 Contribuciones
% ---------------------------------------------------------------------------
% Tres ejes a desarrollar:
%
% - PEDAGÓGICO: trazabilidad real del progreso del estudiante. En lugar de un
%   flag binario "completó / no completó", el educador puede ver QUÉ acción
%   concreta disparó cada target gracias al trigger_event persistido en cada
%   TargetResult, y reconstruir la secuencia del ataque. Esto cambia la
%   pregunta de "¿lo logró?" a "¿cómo lo logró?".
%
% - TÉCNICO Y DE COMUNIDAD: contribución al ecosistema open-source de
%   detección de intrusiones. El código, los patrones de detección
%   (network_payload, process_running con ancestros, sidecar tcpdump por
%   escenario, etc.) y la metodología están en un repositorio público y
%   auditable, frente a las reglas cerradas de los IDS comerciales.
%
% - ARQUITECTÓNICO: separación clara en 4 capas con contratos explícitos.
%   Esto hace al sistema extensible (nuevos tipos de target en Capa 3 sin
%   tocar Capa 1) y portable (cambio de tecnología de virtualización
%   implementando solo EnvironmentProvider). El precedente más cercano,
%   Tectonic, mezcla estas responsabilidades; FlexiPwn las separa.

% ---------------------------------------------------------------------------
% 6.3 Limitaciones reconocidas
% ---------------------------------------------------------------------------
% Buena parte ya está documentada en los capítulos previos y se puede citar
% sin reabrir la discusión:
%
% - GAP DEL POLLING: procesos sub-segundo, fileless, PPID spoofing y técnicas
%   evasivas avanzadas quedan fuera del alcance de detección. Aceptable para
%   el contexto educativo. (Detalle en Capítulo de Implementación, Sección
%   "Limitaciones del Monitoreo por Polling").
%
% - LIMITACIONES DEL NETWORKMONITOR:
%     * Tráfico cifrado (TLS) y protocolos binarios solo exponen metadatos.
%     * Fragmentación de payload entre paquetes TCP — el parser opera por
%       paquete completo, sin reensamblado a nivel de flujo.
%
% - LIMITACIÓN DEL LOGMONITOR: dependencia del flush de la aplicación;
%   exige cooperación del Dockerfile del educador (PYTHONUNBUFFERED=1 o
%   equivalente).
%
% - DEUDA TÉCNICA HEREDADA DE ITERACIONES TEMPRANAS, que queda como
%   limpieza pendiente en el repositorio:
%     * Tipos de target file_exists, http_response_contains y
%       database_query_result definidos en el schema Pydantic pero sin
%       evaluador productor — se eliminarán del schema (no se implementarán).
%     * Campo attacker_ports del EnvironmentConfig sin uso real: el binding
%       del SSH del atacante se asigna automáticamente desde el código.
%     * Dependencia watchdog declarada en pyproject.toml sin uso en código
%       (residuo de la implementación inicial con inotify, sustituida por
%       container.diff()).
%
% - VALIDACIÓN EN INFRAESTRUCTURA REAL: pendiente al momento de redactar.
%   La eficacia del sistema en aula multi-estudiante y la robustez bajo carga
%   están parcialmente cubiertas — ver Capítulo 5.

% ---------------------------------------------------------------------------
% 6.4 Trabajo futuro
% ---------------------------------------------------------------------------
% Organizado en dos horizontes temporales:
%
% CORTO PLAZO — Limpieza y correcciones ya identificadas
%
% - Remover los tipos de target sin productor (file_exists,
%   http_response_contains, database_query_result) del schema Pydantic y de
%   todo material derivado.
% - Remover el campo attacker_ports del EnvironmentConfig.
% - Remover la dependencia watchdog del pyproject.toml.
% - Rehacer el diagrama TikZ de arquitectura para reflejar la interfaz
%   abstracta + tecnología concreta en Capa 1 y agregar el monitor de red en
%   Capa 2.
% - Cerrar la entrega al CLCERT con documentación operativa.
%
% MEDIANO Y LARGO PLAZO — Extensiones del sistema
%
% - PodmanProvider como segundo backend de la interfaz EnvironmentProvider
%   (justificación detallada en Capítulo de Implementación, sección de
%   virtualización). Mejora arquitectónica de seguridad sin afectar capas
%   superiores.
% - INTERFAZ WEB que exponga las mismas operaciones que la CLI clásica
%   manteniendo la lógica de negocio intacta. Útil especialmente para el
%   CLCERT y estudiantes con menos manejo de línea de comandos.
% - DETECCIÓN BASADA EN eBPF / Falco / Sigma rules como complemento del
%   polling, para cerrar el gap de detección de procesos sub-segundo y
%   técnicas evasivas (fileless, PPID spoofing).
% - KATA CONTAINERS o gVisor como runtime alternativo para ejercicios de
%   kernel exploitation, donde el aislamiento de Docker rootless no alcanza.
% - REENSAMBLADO TCP en NetworkMonitor para soportar payloads fragmentados
%   entre paquetes (necesario para escenarios con tráfico de mayor volumen).
% - DOCUMENTACIÓN para creadores de escenarios + plantilla de Dockerfile
%   recomendada, cubriendo el caso del flush y otras buenas prácticas que
%   evitan los puntos ciegos del LogMonitor.

% ---------------------------------------------------------------------------
% 6.5 Reflexión final  (una página corta)
% ---------------------------------------------------------------------------
% - Lugar del proyecto en el ecosistema open-source de cyber ranges
%   educativos en español, y vínculo con el contexto de la universidad y el
%   CLCERT.
% - Aprendizajes personales: lo más valioso de construirlo, lo que se haría
%   distinto si se empezara otra vez.


% ver https://www.overleaf.com/learn/latex/Glossaries
% \input{glosario.tex} % opcional

\nocite{*}
\bibliographystyle{plain}
\bibliography{bibliografia}

% opcional ...
\begin{appendices}
\chapter{Anexo}
\lipsum[50-50]
\end{appendices}
\end{document}
